\documentclass[dvipdfmx,10pt]{jarticle}
\usepackage[dvipdfmx]{graphicx}

\title{ソフトウェア工学チーム開発\break 要件定義書}
\author{6班}
\date{2026年6月19日}

\begin{document}
    \maketitle
    \pagebreak
    \tableofcontents
    \pagebreak
    \section{はじめに}
    本書は、神戸私立工業高等専門学校電子工学科4年次のソフトウェア工学内のチーム開発プロジェクト（６班）の要件定義書である。
    本書では、開発対象システムの要件を明確にし、チーム全体で共通理解を得ることを目的とする。
    
    \section{文書の構成}
    本書のおおよその構成は下記の通りである。
    \begin{itemize}
        \item 第\ref{sec:purpose}節：プロジェクトの目的と概要
        \item 第\ref{sec:overview}節：対象システムの概要
        \item 第\ref{sec:users}節：利用者の分析
        \item 第\ref{sec:usecase}節：システムの利用例
        \item 第\ref{sec:functional}節：機能要件
        \item 第\ref{sec:nonfunctional}節：非機能要件
        \item 第\ref{sec:constraints}節：システム制約
    \end{itemize}

    \pagebreak
    
    \section{プロジェクトの目的と概要}
    \label{sec:purpose}
    \subsection{プロジェクト背景}
    本プロジェクトは、神戸私立工業高等専門学校電子工学科4年次のソフトウェア工学の授業における、チーム開発課題として実施する。
    授業資料では、ToDoアプリを題材として要求分析から成果報告までを行うことが求められており、チームで役割分担をしながら開発を進めることが求められている。

    \subsection{プロジェクトの目的}
    本プロジェクトの目的は、ToDoアプリとして動作するシステムを設計・実装し、課題で示された要求仕様を満たすことである。
    さらに、要求分析、要件定義、設計、実装、テストの流れを通して、ソフトウェア開発の一連の工程を実践的に学ぶことも目的とする。

    \subsection{プロジェクトスコープ}
    本プロジェクトで扱う範囲は、Web アプリとしてのタスク登録、一覧表示、更新、削除、およびステータスによる絞り込みまでとする。
    授業資料で示された最低限の要求仕様を満たしたうえで、開発期間中に実装可能な範囲で追加機能を拡張してよい。
    ただし、開発期間内に完成し、最終提出物として動作確認可能な状態で提出できることを前提とする。
    
    \pagebreak

    \section{対象システムの概要}
    \label{sec:overview}
    \subsection{システムの名称}
    ToDoアプリ

    \subsection{システムの概要}
    今回開発するシステムは、システムの利用者が日々の作業や課題を ToDo として管理できる Web アプリケーションとする。
    利用者はブラウザ上でタスクを登録し、一覧で確認し、内容や状態を更新し、不要になったタスクを削除できる。
    また、タスクをステータスでフィルタリングすることで、未完了の作業や完了済みの作業を見分けやすくする。

    \subsection{開発環境}
    本システムは下記の開発環境で実装するWebアプリケーションとする。
    \begin{itemize}
        \item 技術スタック：HTML, CSS
        \item 使用ツール：各自の開発環境、テスト環境、必要に応じて設計支援ツール
        \item プラットフォーム：Web ブラウザ
    \end{itemize}
\pagebreak
    \section{利用者の分析}
    \label{sec:users}
    \subsection{ユーザー層の分類}
    想定利用者は、日常的に複数の課題や締切を管理したい学生である。
    特に、授業課題、提出物、レポート、試験準備などの期限付きタスクを扱う利用者を主な対象とする。

    \subsection{ユーザーの特性}
    利用者は、どのタスクから手を付けるべきかを素早く判断したい傾向がある。
    そのため、入力項目が明確で、期限や重要度、状態を一覧から把握しやすいことが重要となる。
    特に、Web アプリとして直感的に操作できる画面構成が必要となる。

    \pagebreak

    \section{利用者のニーズと課題}
    \label{sec:needs}
    \subsection{現在の問題点}
    紙のメモや分散した記録では、タスクの抜け漏れや重複管理が発生しやすい。
    また、何を先にやるべきかが分からないまま課題が増えると、今やるべきタスクを見失いやすい。

    \subsection{ユーザーの要望}
    利用者は、タスクを簡単に登録でき、あとから内容を編集でき、不要なものを削除できることを求める。
    さらに、学生向けに使いやすく、期限が近いタスクを見逃しにくいことや、Google Classroomのように締切を意識できる機能を求める。

    \subsection{解決すべき課題}
    本システムは、タスクの登録から整理までを一つの画面または一連の操作で行えるようにし、管理の手間を減らす必要がある。
    加えて、入力値の不備を防ぎ、誤った内容の登録を避けることで、信頼性の高い管理を実現する。
    どのタスクを優先すべきかが分かるように、期限・重要度・ステータス・通知を組み合わせて表示する必要がある。

    \pagebreak

    \section{システムの利用例（ユースケース）}
    \label{sec:usecase}
    \subsection{ユースケース図}
    利用者は、タスク登録、一覧確認、更新、削除、ステータスによるフィルタリングを行う。
    これらに加えて、期限が近いタスクを確認し、通知によって締切を意識することが主要なユースケースとなる。

    \subsection{主要なユースケース}
    \begin{itemize}
        \item 利用者がタスク名、ジャンル、期限、メモ、重要度、ステータスを入力して新しいタスクを登録する。
        \item 利用者が登録済みタスクの一覧を確認し、今日・今週・来週締切のタスクを優先的に把握する。
        \item 利用者がタスクのステータスを一覧上のプルダウンで変更する。
        \item 利用者が詳細編集画面からタスク内容を更新する。
        \item 利用者が詳細編集画面から不要になったタスクを削除する。
        \item 利用者が完了・未完了などのステータスで表示対象を絞り込む。
    \end{itemize}

    \pagebreak

    \section{機能要件}
    \label{sec:functional}
    \subsection{機能一覧}
    本システムに実装する機能は下記の通りである。
    \begin{itemize}
        \item タスク登録機能
        \item タスク一覧表示機能
        \item タスク更新機能
        \item タスク削除機能
        \item ステータスフィルタ機能
        \item 期限別一覧表示機能
        \item 通知・リマインド機能
        \item 入力値検証機能
    \end{itemize}

    \subsection{詳細な機能説明}
    \begin{itemize}
        \item タスク登録機能では、利用者がタスクを新規追加できる。登録項目はタイトル、ジャンル、期限、メモ、重要度、ステータスとする。
        \item タスク一覧表示機能では、登録済みのタスクを一覧で確認できる。初期画面では今日・今週・来週締切のタスクを優先表示し、それ以降は折りたたみ等で確認できるようにする。
        \item タスク更新機能では、タスク名やステータスなどの内容を編集できる。タスク一覧の各項目にあるプルダウンでステータス変更を行い、詳細編集画面から内容変更を行えるようにする。
        \item タスク削除機能では、不要になったタスクを詳細編集画面などから削除できる。
        \item ステータスフィルタ機能では、未着手、進行中、提出前、完了済みなどの状態ごとに表示対象を切り替えられる。
        \item 期限別一覧表示機能では、締切の近いタスクを優先して確認できるようにする。
        \item 通知・リマインド機能では、期限が近いタスクを利用者に知らせ、締切忘れを防ぐ。これは初期搭載機能とする。
        \item 入力値検証機能では、空文字の登録を禁止し、140字以上の入力を受け付けないなど、仕様に反する入力を防ぐ。
    \end{itemize}

    \section{非機能要件}
    \label{sec:nonfunctional}
    \subsection{パフォーマンス要件}
    タスクの登録、更新、一覧表示、フィルタリングは、利用者が待ち時間を強く意識しない程度の応答速度で動作することを目標とする。

    \subsection{セキュリティ要件}
    入力内容の検証を行い、不正な値や空入力による不整合を防ぐ。
    また、利用者が意図しないデータ破壊が起こらないよう、削除や更新の操作は明示的に行えるようにする。

    \subsection{ユーザビリティ要件}
    Web ブラウザ上で初めて使う利用者が操作手順を理解しやすい構成とする。
    タスクの状態が一覧から把握しやすく、操作方法が過度に複雑にならないことを重視する。

    \subsection{その他の非機能要件}
    授業要件として、テストを実施し、仕様に対して正しく動作することを確認する。
    実装は HTML, CSS を用いた Web アプリとし、開発期間内に完成できる規模に収める必要がある。
    また、学生が日常的に使いやすいことを重視し、登録・確認・更新の流れが直感的に理解できることを目指す。

    \pagebreak

    \section{システム制約}
    \label{sec:constraints}
    \subsection{技術的制約}
    実装は HTML, CSS を用いた Web アプリとする。
    利用環境は Web ブラウザを前提とし、授業の要求仕様を満たす必要がある。
    ただし、学生向けの利用を前提とし、Google Classroomのような期限リマインド機能を初期搭載する。

    \subsection{スケジュール制約}
    授業資料のスケジュールでは、中間テスト後に要求分析、要件定義、DFD、クラス図、開発、成果報告会、納品へと進む。
    要件定義書は中間段階で一旦提出し、その後の設計・開発に反映する必要がある。

    \subsection{予算制約}
    特別な予算を前提としない。
    基本的には各自の開発環境を用いて実施し、追加の有償サービスや機材を前提としない。

    \pagebreak
    
    \section{参考資料}
    \begin{itemize}
        \item ソフトウェア工学 第7回 チーム開発 配布資料
        \item 要求仕様：ToDoアプリ
    \end{itemize}

\end{document}